%-*-latex-*-
%-*-latex-*-
\newcommand\COURSE{ciss350}
\newcommand\ASSESSMENT{q10802}
\newcommand\ASSESSMENTTYPE{Quiz}
\newcommand\POINTS{\textwhite{xxx/xxx}}

\input{myquizpreamble}
\input{yliow}
\input{\COURSE}
%-*-latex-*-

%-*-latex-*-

%-*-latex-*-
\renewcommand\TITLE{\ASSESSMENTTYPE \ \ASSESSMENT}

\renewcommand\EMAIL{}
\renewcommand\AUTHOR{eagraham1@cougars.ccis.edu}

\textwidth=6in
\begin{document}
\topmatterthree

You are given this class for the row of a hashtable:
\begin{console}[fontsize=\footnotesize]
const int Available = 0;
const int NotAvailable = 1;
const int Deleted = 2;

class HashtableRow
{
public:
    HashtableRow(int flag, int stu_id, double height)
        : flag_(flag), stu_id_(std_id), height_(height)
    {}
    int flag_;
    int stu_id_;
    double height_;
};

int hash(int x)
{
    return 2 * x + 3;
}
\end{console}
Complete the following function to insert a row into the given hashtable.
This uses non-separate chaining.
The index of the row where the insert occurs is returned.
If the table is full, -1 is returned.
$n$ is the size of the hashtable.
\\
\ANSWER
\begin{answercode}
int hashtable_insert(HashtableRow hashtable[], int n, int key, double value)
{
  int index = -1;
  int h = hash(key) % n;
  int start = h;

  while(true)
  {
    if(hashtable[h].flag == 2) // Deleted
      index = h
    else if(hashtable[h].flag == 0) // Available
    {
      if(index == -1)
        index = h;
        
      hashtable[h].stu_id_ = key;
      hashtable[h].height_ = value;
      hashtable[h].flag = 1;
      
      return index;
    }
    else if(hashtable[h].flag == 1) // Not Available
    {
      if(hashtable[h].stu_id_ == key)
        return -1;
    }

    h = (h + 1) % n;

    if(h == start)
      return -1;
  }
}
\end{answercode}

%------------------------------------------------------------------------------
%-*-latex-*-
\newpage
\input{instructions}
\end{document}

